\documentclass[11pt]{amsart}
\usepackage[T1]{fontenc}
\usepackage{amsmath,amssymb,amsthm}
\usepackage{booktabs}
\usepackage[hidelinks]{hyperref}
\hypersetup{pdfauthor={The UNICO/NOUS pipeline},
  pdftitle={A kernel-certified core for the study of Agrawal's conjecture at r=5}}

\newtheorem{theorem}{Theorem}[section]
\newtheorem{lemma}[theorem]{Lemma}
\newtheorem{corollary}[theorem]{Corollary}
\newtheorem{conjecture}[theorem]{Conjecture}
\theoremstyle{remark}
\newtheorem{remark}[theorem]{Remark}

\newcommand{\eps}{\varepsilon}
\newcommand{\Z}{\mathbf{Z}}
\newcommand{\F}{\mathbf{F}}
\newcommand{\lean}[1]{\texttt{#1}}
\emergencystretch=1.5em

\begin{document}

\title[A kernel-certified core at \(r=5\)]{A kernel-certified core for
the study of Agrawal's conjecture at \(r=5\)}

\author{The UNICO/NOUS pipeline}
\thanks{An autonomous multi-model pipeline (Claude Fable~5, Claude
Opus~4.8, GPT-5.6-Sol, with an independent external review pass),
orchestrated by Daniele Cappello, who carried the messages, audited
the artifacts and made no mathematical contribution. Repository with
the Lean sources and all computational certificates:
\url{https://github.com/Solarys431/agrawal-r5}.}

\begin{abstract}
This note collects the part of our structural study of Agrawal's
conjecture at \(r=5\) that is mechanically verified. Every statement
called a theorem below is proved in Lean~4 over Mathlib, with no
\lean{sorry}, no \lean{admit} and no added axioms; the corresponding
declaration name is given with each statement. Every computational
claim carries a replayable certificate with its detection criterion
embedded and SHA-256 manifests. We then pose, without any claim of
proof, the two open problems that our working material identifies as
central. Nothing in this note asserts the conjecture.
\end{abstract}

\maketitle

\section{Setting}

Agrawal's conjecture (arising from the AKS primality test
\cite{AKS}; see also \cite{LPremarks,LPgauss,Vana}) states: if \(r\) is prime,
\(r\nmid n\), and
\((X-1)^n\equiv X^n-1 \pmod{n,\,X^r-1}\),
then \(n\) is prime or \(n^2\equiv1\pmod r\). Our study concerns
\(r=5\). This note isolates its mechanically certified core; the
statements are elementary, and their interest lies in the role they
play in the wider study and in the fact that a proof assistant has
checked every one of them.

In Sections 2--4, \(p\) is an odd prime; in Section~5, \(q\) denotes
an arbitrary prime factor of \(n\), with \(q=2\) allowed. Let
\[
R_p \;=\; (\Z/p)[X]\,/\,(X^2-X-1),
\qquad
\eps \;=\; \text{the class of } X,
\]
so that \(\eps^2=\eps+1\), and set \(\sqrt5:=2\eps-1\), which
satisfies \((\sqrt5)^2=5\) in \(R_p\)
(\lean{GoldenRing}, \lean{eps\_sq}, \lean{sqrt5\_sq}).
Let \(F_n\), \(L_n\) denote the Fibonacci and Lucas numbers, with
\(L_n=F_{n+1}+F_{n-1}\) (\lean{lucas}). Define
\[
A_n=\begin{cases}L_n,& n \text{ even},\\ F_n,& n\text{ odd},\end{cases}
\qquad
H_n=\gcd\bigl(A_n,\,5^{\,n-1}+1\bigr),
\]
and, dually, \(B_n=F_n\) (\(n\) even), \(L_n\) (\(n\) odd), with
\(J_n=\gcd\bigl(B_n,\,5^{\,n-1}-1\bigr)\).

\section{The golden Frobenius and the inertia of \texorpdfstring{\(J_n\)}{Jn}}

\begin{theorem}[\lean{golden\_frobenius}]\label{thm:frob}
Assume \(5\) is not a square in \(\Z/p\) and \(p\neq2\). Then
\(\eps^p=1-\eps\) in \(R_p\).
\end{theorem}

\begin{proof}[Certified proof sketch]
By Euler's criterion (\lean{euler\_nonsquare}) we have
\(5^{(p-1)/2}=-1\) in \(\Z/p\), hence
\[
(\sqrt5)^p=(5)^{(p-1)/2}\sqrt5=-\sqrt5
\qquad(\lean{sqrt5\_pow\_p}).
\]
Expanding \((2\eps-1)^p\) by the Frobenius endomorphism and
cancelling the unit \(2\) gives the claim.
\end{proof}

\begin{corollary}[\lean{golden\_pow\_p\_succ}]\label{cor:halfperiod}
Under the same hypotheses, \(\eps^{\,p+1}=-1\) in \(R_p\).
\end{corollary}

\begin{theorem}[\lean{inertia\_J}]\label{thm:inertia}
Assume \(5\) is not a square in \(\Z/p\) and \(p\neq2\). Then for no
\(n\ge1\) do \(\eps^{2n}=1\) in \(R_p\) and \(5^{\,n-1}=1\) in
\(\Z/p\) hold simultaneously.
\end{theorem}

\begin{proof}[Certified proof sketch]
If \(n\) is odd, raising \(\eps^{2n}=1\) to the power \((p+1)/2\) and
comparing with \((\eps^{p+1})^n=(-1)^n=-1\) gives \(1=-1\), absurd for
\(p\neq2\). If \(n\) is even, \(5=5^{\,n}=(5^{\,n/2})^2\) exhibits
\(5\) as a square, against the hypothesis.
\end{proof}

Via the identity \(\eps^{\,n+1}=F_{n+1}\eps+F_n\) in \(R_p\)
(\lean{golden\_pow\_fib} and its predecessor form
\lean{golden\_pow\_pred}) one obtains the divisibility forms, in
which the hypothesis is purely arithmetic by quadratic reciprocity
(\lean{not\_isSquare\_five}):

\begin{theorem}[\lean{inertia\_J\_fib\_mod5},
\lean{inertia\_J\_lucas\_mod5}]\label{thm:inertiadvd}
Let \(p\equiv\pm2\pmod5\), \(p\neq2\). For even \(n\ge2\), \(p\) does
not divide \(\gcd(F_n,\,5^{\,n-1}-1)\); for odd \(n\ge1\), \(p\) does
not divide \(\gcd(L_n,\,5^{\,n-1}-1)\). Equivalently, no prime inert
in \(\mathbf{Q}(\sqrt5)\) divides any \(J_n\).
\end{theorem}

\section{The canonical support witness}

\begin{theorem}[\lean{support\_witness}]\label{thm:support}
Let \(p\equiv2\pmod5\), \(p\neq2\), and put \(n_0=(p+1)/2\). Then
\(n_0\equiv4\pmod5\), \(p\mid 5^{\,n_0-1}+1\), and \(p\) divides
\(L_{n_0}\) if \(n_0\) is even, \(F_{n_0}\) if \(n_0\) is odd. In
particular \(p\mid H_{n_0}\), with an index in the class
\(4\bmod5\).
\end{theorem}

\section{The golden--cyclotomic entanglement}

\begin{theorem}[\lean{entanglement}]\label{thm:ent}
In any commutative ring, if
\(\zeta^4+\zeta^3+\zeta^2+\zeta+1=0\), then
\[
\zeta^3\,(1+\zeta+\zeta^4)^2\,(\zeta-1)^4 \;=\; 5 .
\]
In particular this holds in \(\Z[\zeta_5]\) with
\(\eps=1+\zeta+\zeta^4\) the image of the golden unit.
\end{theorem}

The proof is a single \lean{linear\_combination} with an explicit
cofactor; \(\zeta^5=1\) follows from the same hypothesis
(\lean{zeta\_pow\_five}).

\section{The Fermat-shadow glue}

\begin{lemma}[\lean{mul\_dvd\_gcd\_mul}]\label{lem:index}
If \(x\mid M\) and \(y\mid M\), then \(xy\mid\gcd(x,y)\,M\).
\end{lemma}

\begin{theorem}[\lean{fermat\_shadow}]\label{thm:shadow}
Let \(n\) be squarefree with \(5\nmid n\), and suppose that for every
prime \(q\mid n\)
\[
5^{\,n/q-1}\equiv1\pmod q .
\tag{\(\ast\)}
\]
Then \(n\mid 5^{\,n-1}-1\); that is, \(n\) is a Fermat pseudoprime to
base \(5\).
\end{theorem}

\begin{remark}
Hypothesis \((\ast)\) is the local norm constraint of the wider study;
in Theorem~\ref{thm:shadow} it is an explicit hypothesis. The next
theorem removes it: the Agrawal congruence itself forces base-\(5\)
pseudoprimality, unconditionally.
\end{remark}

\begin{theorem}[\lean{agrawal\_fermat\_shadow}]\label{thm:bridge}
Let \(n\) be squarefree with \(5\nmid n\), and suppose the Agrawal
congruence at \(r=5\) holds:
\((X-1)^n\equiv X^n-1\pmod{n,\,X^5-1}\).
Then \(n\mid 5^{\,n-1}-1\). In particular every squarefree
counterexample to Agrawal's conjecture at \(r=5\) prime to \(10\) is
a Fermat pseudoprime to base \(5\).
\end{theorem}

\begin{proof}[Certified proof sketch]
Reduce the congruence mod a prime \(q\mid n\); it becomes a
polynomial identity in \((\Z/q)[X]\), which can be evaluated at the
four elements \(z^j\) (\(j=1,\dots,4\)) of the quotient ring
\((\Z/q)[X]/\Phi_5\) (not a field in general), where \(z^5=1\) kills
the modulus. Multiplying the four evaluated identities and using the
explicit-cofactor identity \((z-1)(z^2-1)(z^3-1)(z^4-1)=5\)
(\lean{prod\_pow\_sub\_one}) gives \(5^{\,n}=5\) in the quotient
ring; injectivity of the constant embedding descends this equality
to \(\Z/q\), hence \(5^{\,n-1}=1\) there (\lean{agrawal\_local}),
and the squarefree gluing finishes. No Frobenius descent and no norm
maps are needed.
\end{proof}

\section{Computational certificates}

Independently of the kernel-checked statements above, the repository
carries replayable certificates, each embedding its own detection
criterion and reconstruction data, with SHA-256 manifests:

\begin{enumerate}
\item seven \emph{certified-empty fibres} of a three-factor sieve
associated with the study (pairs
\((13,37)\), \((17,113)\), \((13,277)\), \((23,167)\), \((7,823)\),
\((83,347)\), \((463,1327)\)): for each pair, the finitely many
levels required by the certificate are integrally factored, all prime
factors are proven prime, and no prime in the certificate's detector
classes survives;
\item a self-certifying census of \(H_n\) for all
\(n\equiv4\pmod5\), \(n\le10^5\): \(9{,}725\) nontrivial values,
all factorizations proven, and \emph{no} prime factor
\(\equiv\pm1\pmod5\) occurs.
\end{enumerate}

These are certificates of finite computations; they prove exactly
what they state and nothing beyond, and both replay from the
repository alone. Separately, the repository freezes as
\emph{hash-only provenance} (hashes, not artifacts) a SHA-256 index
of the working material of the study, including a prime-by-prime
corpus up to \(10^9\); nothing in this note relies on it.

\section{Two open problems}

We pose the two statements that our working material identifies as
the central open problems of the \(r=5\) study. We claim no proof and
no partial proof here.

\begin{conjecture}[golden inertia]\label{conj:inertia}
For every \(n\equiv4\pmod5\), every odd prime factor of \(H_n\)
distinct from \(5\) is \(\equiv\pm2\pmod5\).
\end{conjecture}

Theorem~\ref{thm:support} shows that every prime
\(p\equiv2\pmod5\) does occur in some \(H_{n}\) with
\(n\equiv4\pmod5\); the census in item (2) above verifies
Conjecture~\ref{conj:inertia} for all \(n\le10^5\).

\begin{conjecture}[fibre emptiness]\label{conj:fibres}
The pattern of item (1) above is general: for every admissible pair,
the corresponding fibre is empty.
\end{conjecture}

\section{Formalization data}

Nine modules, \(834\) source lines, \(56\) named declarations; no
\lean{sorry}, no \lean{admit}, no axioms beyond Mathlib's. Mathlib is
pinned to a fixed upstream commit; a clean clone builds with
\lean{lake exe cache get \&\& lake build}. The table mapping each
statement of this note to its declaration and file is in the
repository \lean{README}.

\begin{thebibliography}{9}

\bibitem{AKS}
M.~Agrawal, N.~Kayal, and N.~Saxena,
\emph{PRIMES is in P},
Ann. of Math. 160 (2004), 781--793.

\bibitem{LPremarks}
H.~W. Lenstra, Jr. and C.~Pomerance,
\emph{Remarks on Agrawal's conjecture},
AIM note, 2003.

\bibitem{LPgauss}
H.~W. Lenstra, Jr. and C.~Pomerance,
\emph{Primality testing with Gaussian periods},
J. Eur. Math. Soc. 21 (2019), 1229--1269.

\bibitem{Vana}
T.~V\'a\v na,
\emph{Agrawal's conjecture and Carmichael numbers},
Student Scientific Conference, Comenius University Bratislava, 2009.

\end{thebibliography}

\end{document}
